% 原文 pp.13--19、およびp.20冒頭の継続段落。本文を段落単位で照合。
\sectionlink{デコヒーレンスの基本的な形式と解釈}
% 原文 p.13
\hypertarget{srcpage:13}{}\label{srcpage:13}
本章の前半では、デコヒーレンスの研究計画を支える量子力学の形式と解釈のいくつかを紹介する。まず量子状態を論じる（2.1節）。量子力学が古典物理学と異なる点の多くは、量子状態の概念と解釈を中心に整理できる。古典物理学では「系の状態」は、位置、運動量、温度などの物理量の値を列挙したものを意味するにすぎない。量子物理学では、系は抽象的な状態ベクトルによって記述され、一般にこれを観測される物理的性質と結び付けるには、量子測定という付加的な概念を必要とする。しかも、この概念はよくても曖昧にしか定義されていない。

続く2.2節では量子力学の重ね合わせ原理を論じる。抽象的な量子状態という概念が古典力学と量子力学の最初の重要な相違であったとすれば、重ね合わせ原理は両者の隔たりをさらに広げる。デコヒーレンスの大部分は量子的重ね合わせの問題を中心にしている。すなわち、自然界でどの重ね合わせを観測できるかを教える中心的な機構であり、重ね合わせ原理によって原理上は許される重ね合わせの大半が、実際には周囲の世界で見られない理由を説明する。

2.3節では、デコヒーレンスにとって重要な、もう一つの量子力学特有の概念、すなわち量子もつれを導入する。量子もつれは、量子計算などの技術的応用において量子力学の力を利用するための中心的な要素となっている。また、デコヒーレンスの基礎にある機構でもある。2.4節では、デコヒーレンスの形式的記述に重要な道具となる密度行列の形式と解釈を論じる。2.5節では量子力学の測定問題と、それに関連する量子から古典への移行の問題をやや詳しく紹介する。デコヒーレンスはこれらの問題と密接に関係し、本章および第8章で見るように、少なくとも部分的な解決を与える。

2.6節では二重スリット実験を例として、デコヒーレンスに不可欠な量子もつれの帰結、すなわち、ある系ともつれた相手に「どちらの経路を通ったか」という情報が符号化されることを論じる。この説明によって、系と環境の量子もつれを介した、環境による系の「監視」の帰結としてデコヒーレンスを理解する道筋が開かれる。

% 原文 p.14
\hypertarget{srcpage:14}{}\label{srcpage:14}
続く2.7、2.8、2.9節では、この環境による監視の三つの主要な帰結を詳しく論じる。すなわち、系の水準での干渉効果の抑制、環境との量子もつれに最も影響されにくい準古典的な優先状態の選択、そして、これらの優先状態に関する情報が環境内に頑健かつ冗長に符号化されることである。第1章ですでに指摘したように、狭義のデコヒーレンスはしばしば第一の帰結（干渉の減衰）のみを指し、第二の帰結（優先状態の選択）は「環境誘起超選択」という別の用語で呼ばれる。しかし両現象は起源を共有し、不可分に結び付いている。干渉の減衰は、干渉が抑制される対象となる特定の状態の集合について述べる必要があるからである。この広い見方では、「デコヒーレンスの研究計画」は、環境による監視の上述の三つの帰結をすべて含むものと考えるべきである。

2.10節では、具体的で単純なデコヒーレンス模型を通じてこれらの考えを検討する。残りの各節では、いくつかの重要な追加の話題を論じ、全体像を完成させる。デコヒーレンスと散逸（2.11節）、およびデコヒーレンスと古典的雑音（2.12節）の根本的な相違を強調する。2.13節ではデコヒーレンスが可逆となりうる例を述べ、関連する量子的「消去」の過程を分析する。2.14節では、デコヒーレンスの研究計画の重要な概念的問題である、宇宙を部分系に分解する問題を論じる。最後に2.15節では、さらにいくつかの形式的な道具、すなわちSchmidt分解定理（2.15.1節）、Wigner表示（2.15.2節）、非純粋状態の純粋化（2.15.3節）、およびKrausの演算子和形式（2.15.4節）を概説する。

\subsectionlink{量子状態の概念と解釈}
\subsubsectionlink{古典状態と量子状態}
古典物理学における物理系の「状態」という概念は、きわめて直観的である。通常、たとえば運動する物体の位置や運動量のような、関心のある測定可能な量に注目し、それらの量に数学的記号（「$x$」や「$p$」など）を割り当てる。すると、これらの記号に数値を与えるだけで物体の運動状態が指定される。言い換えれば、物体の物理的性質（したがって物理世界の実在物）と、理論におけるその形式的・数学的表現との間に一対一の対応が存在する。

% 原文 p.15
\hypertarget{srcpage:15}{}\label{srcpage:15}
もちろん古典物理学でも、ここで用いたNewton力学の例ほど容易には、この直接の対応が確立されない場合を考えることはできる。たとえば歴史的には、熱力学理論における温度の形式的定義を、物理的な温度概念の基礎となる分子過程に結び付けるのは、かなり難しいことが判明した。しかし通常は、他の物理量や現象を参照することによって、少なくともある水準ではこの同定の問題を解決できた。たとえば、集合的な熱力学変数を、物質の集まりに含まれる微視的構成要素の振る舞いに結び付けるのである。

20世紀初頭の量子論の登場によって、物理世界と理論内の数学的表現との、この単純な全単射的対応は突然終わりを迎えた。経験に現れる「客観的に存在する」物理的性質に直接対応する直観的な記号によって物理系の状態を記述する代わりに、量子力学で利用できるのは、同じく抽象的なHilbertベクトル空間内のベクトル（より一般には射線）として定義される、抽象的な量子状態だけである。

この抽象化に伴う概念的飛躍の大きさは、いくら強調しても足りない。実際、量子論の初期以来、数え切れない物理学者や哲学者が取り組んできた「量子力学の解釈」をめぐる議論は、その大部分がまさに、抽象的な量子状態を「外にある物理的実在」とどう関係付けるかという問いに根差している。

教科書的な量子力学によれば、経験上なじみのある物理量とのつながりは、間接的にのみ作られる。すなわち物理量の測定を通じてであり、その物理量、つまり観測可能量（「現実世界」の側）は、Hilbert空間内のHermite演算子（「形式的理論」の側）によって表される。通常用いられる収縮（または射影）公準によれば、測定は量子状態を、測定された観測可能量を表す演算子の固有状態の一つへと瞬間的に変化させる。それぞれの固有状態が得られる確率はBorn則\cite{ref019}で与えられる（原文35頁の脚注9を参照）。

測定される演算子の固有状態は、対応する物理量が測定で取りうる、互いに異なる確定した「値」を表す。したがって、これらの固有状態に結び付いた概念は、上で論じた古典物理学のものに似ている（もちろん、測定する観測可能量は、スピンの場合のように、明確に非古典的でありうる）。それゆえ測定は、ある程度まで、数学的形式と系の「客観的に存在する物理的性質」との一対一の対応、つまり古典物理学でなじみ深い概念への回帰を可能にする。

しかしここには、量子力学特有のもう一つの注意点が潜んでいる。多くの観測可能量が互いに両立しないため（形式的には、対応する演算子の非可換性によって表現される）、量子状態は一般に、きわめて少数の演算子観測可能量の同時固有状態でしかない。したがって量子系に帰属させられる確定した物理的性質の数は限られ、測定を追加すると一般に系の状態は変化する（「乱される」）。もちろん、まったくの幸運や事前知識によって、測定前の系の量子状態にたまたま一致する固有状態を持つ演算子観測可能量を測定する場合は別である。

% 原文 p.16
\hypertarget{srcpage:16}{}\label{srcpage:16}
量子状態のこの壊れやすさは、有名な複製不可能定理（7.4.1節）にも反映されている。この定理によれば、未知の量子状態を複製することは一般に不可能である\cite{ref020,ref021}。これを別の言い方で表すと、個々の系の未知の量子状態を、その系だけに対する測定によって一意に決定することは不可能である\cite{ref022,ref023,ref024,ref025,ref026}。

言うまでもなく、この状況は古典物理学と際立った対照をなす。古典物理学では、さらに別の物理量を、任意の順序でいくつでも測定することにより、系の物理的性質の「一覧」を拡張し、したがって状態をより完全に指定できる。さらに、多数の独立した観測者が、初めは系の状態についてまったく知らなくても、状態を乱す危険にさらされることなく、そのような測定を行い、結果について一致できる。もちろん、このように古典状態が客観的にあらかじめ存在することが、私たちの深く根付いた「古典的実在」の概念をもたらした。

以下では、曖昧さが生じる場合を除き、物理的観測可能量を表すHermite演算子を、より正確な「演算子観測可能量」という語の代わりに、単に「観測可能量」と呼ぶことが多い。それでも、形式的演算子と物理量との同一視は量子論に公理的に導入されるものであり、自明ではないことを念頭に置いてほしい。ただし、観測可能量の概念を創発的なものと見なし、量子状態の形式だけから導くこともできる。この見方はJohn Bellの「beable」の探究\cite{ref027}の精神に沿うものである（その例は文献\cite{ref017}の2.2節を参照）。物理量の測定を、過度に単純化して「演算子の測定」と記述することにより、このような演算子について素朴実在論に陥る危険は、すでに指摘されている\cite{ref028}。

\subsubsectionlink{量子状態の確率的性格}
前節で論じた量子状態の性質を踏まえ、これらの状態は、観測されるさまざまな「古典」状態への「可能性」だけを表すのだと、しばしば主張されてきた。しかし同時に、現在の知識によれば、量子状態は量子系の完全な記述を表すことを強調するのが重要である。すなわち、量子状態は、系の物理状態について言いうるすべてを包含している。それでも一般には、量子状態は、測定で具体的にどの結果が得られるかを教えず、可能なさまざまな結果の確率だけを教える。量子力学の、この一見内在的な確率的性格は、この理論を古典物理学から区別する中心的な特徴の一つである。実験では、この確率的側面は、まったく同じ量子状態に準備した系の集まりについて同じ物理量を測定すると、一般に異なる結果の集合が得られることとして現れる。

% 原文 p.17
\hypertarget{srcpage:17}{}\label{srcpage:17}
よく知られているように、Einsteinは生涯を通じて、自然のこの見かけ上内在的な偶然性を受け入れることに消極的だった\cite{ref029}。それは「神は宇宙でさいころを振らない」という有名な言葉に表れている。この不満は他の人々にも共有され、彼らは、量子状態だけでは系の物理状態を完全には指定できないという意味で、量子力学は完全な理論ではないと提案して、逃れる道を求めた。その場合、「同一に準備された系」の集団は、実際には物理的に同一な系の集まりを表していないことになる。集まりの各系は同じ量子状態で記述されても、それぞれの個々の系に対する、可能なすべての測定の結果を決める「完全な状態」は、集団の各要素で同一ではないというのである。

そのような手法は通常、隠れた変数理論と呼ばれる。これらは、系に対して行われる任意の測定の具体的結果を、その値が決定する「隠れた変数」を仮定して導入し、量子状態を補うことにより、根本的な水準で決定論を回復しようとする。したがって、量子状態と、関係する隠れた変数の値との組合せだけが、系の完全な物理状態を指定する。すべて同じ量子状態に準備された系の集団は、隠れた変数の異なる値の範囲によって記述されることになる。こうすれば、測定結果が異なることを完全に決定論的に説明できる。量子力学が私たちに確率的に見えるのは、隠れた変数の具体的な値を知る手段がないからにすぎない（それゆえ「隠れた」と呼ばれる）。したがって量子力学の確率は、純粋に認識論的なものになる。古典統計力学と同じく、確率は私たちの主観的な知識の不足を表すにすぎず、物理学の根本水準における客観的な非決定性を表すものではない。

しかし実際には、実験によって十分に確かめられた標準量子力学の予言に反しない、隠れた変数の手法を構築するのは、きわめて難しいことが判明している。最も有名なのは、John Bellが1964年の画期的な論文で提示した不等式であり、あらゆる局所的な隠れた変数理論はこれを満たすが、標準量子力学はこれを破る\cite{ref030}（後にBellがさらに洗練した説明\cite{ref027,ref031}も参照）。多数の実験がBellの不等式の明確な破れを示し、局所的な隠れた変数理論を、ほぼ決定的に排除した\footnote{実験結果の解釈には、ごくわずかで、しかも減少し続けてはいるが、なお余地が残っている。主な「抜け穴」は、検出効率が100\%に満たない不完全な検出器によるものである\cite{ref032}。もう一つの問題は「通信の抜け穴」と呼ばれるもので、現在ではWeihsらの実験\cite{ref033}によって閉じられたと一般に考えられている。}。これらの結果は、隠れた変数理論が機能するためには、高度に非局所的である必要を示している。そのような非局所的な隠れた変数理論の代表として、おそらく最も有名なのはde Broglie--Bohmの手法\cite{ref034,ref035,ref036,ref037}である（8.5節）。ここでは位置を隠れた変数に選び、系の各粒子は常に確定した位置を持つ。

% 原文 p.18
\hypertarget{srcpage:18}{}\label{srcpage:18}
さらにKochenとSpecker\cite{ref038}は、隠れた変数理論は必ず文脈依存でなければならないことを示した。すなわち、隠れた変数によって量子系に帰属させる物理量の具体的な値は、一般に測定の文脈、つまり最終的にその値をどのような仕方で測定するかに依存する。たとえば、その系について他のどの観測可能量を同時に測定するかによって、系は、ある観測可能量の異なる値を持つことになる。したがって文脈依存性は、そもそも隠れた変数の研究計画の動機となった中心的な考えを放棄するよう迫る。その考えとは、物理世界はそれに対して行うどんな測定からも独立であり、測定は、隠れた変数の値によって決まる、あらかじめ存在する物理量の値を単に明らかにするのだというものである。このように、いかなる隠れた変数の概念を仮定しても、提案どおりに機能するには、精緻で、高度に非局所的かつ文脈依存の機構が必要になることは明らかである。

\subsubsectionlink{量子状態の存在論的地位}
もう一つの論争は、量子状態の存在論的地位という問いを中心に展開してきた（簡潔で読みやすい歴史的概説は文献\cite{ref039}を参照）。量子力学の初期に、Schrödingerは実空間内の狭い波束を、実際の物理的粒子と同一視しようとした（第3章の冒頭）。この方策には二つの中心的な問題があった。第一に、初めに局在した波束は一般に、空間の広い領域へ急速に広がる傾向があり、定義上空間に局在する粒子という概念とは両立しない。第3章で見るように、この広がりは環境との相互作用によって実効的に抑制されるが、もちろん当時のSchrödingerはその過程を知らなかった。第二に、3次元実空間内の$N>1$個の粒子の量子状態を記述する波動関数は、$3N$次元のHilbert空間に存在し、もはや経験上なじみのある3次元空間内には存在しない。

その後、BornとPauliは別の手法を採り、量子状態は確率振幅を表す、すなわち、系に対して行いうるすべての測定の結果の確率を指定する、という有名な解釈を定式化した。そのうえで測定という行為が、これらの潜在的性質を動的に「現実化する」という根本的役割を果たすと仮定した。Wheeler\cite[p.~182]{ref040}（文献\cite{ref041}も参照）は、この見解を「どんな素過程［すなわち量子］現象も、それが記録された（観測された）現象になるまでは、現象ではない」と端的に表した。Heisenbergの「粒子の軌道は、私たちがそれを観測するという行為によって作られる」\footnote{ドイツ語の原文は「Die Bahn entsteht erst dadurch, daß wir sie beobachten.」である。}\cite[p.~185]{ref042}という言葉や、Bornに宛てたPauliの手紙にある、「観測の際に確定した位置$x_0$が現れること（中略）は、自然法則の外にある創造と見なされる」\footnote{原文は「Das Erscheinen eines bestimmten Ortes $x_0$ bei der Beobachtung (.\, .\, .) wird dann als außerhalb der Naturgesetze stehende Schöpfung aufgefaßt.」である。}\cite{ref043}という示唆にも、同様の考えが見られる。このように、広く受け入れられている波動関数のBorn--Pauli解釈は、量子状態が物理的実在を直接表すと見なす意味で、純粋に存在論的なものでもなく、私たちの主観的知識の不足だけを表すという意味で、単純に認識論的なものでもない。

% 原文 p.19
\hypertarget{srcpage:19}{}\label{srcpage:19}
しかしこの解釈は、「観測による現実化」の存在論的役割とは何か、そしてこの過程を力学的な言葉でどう物理的に説明するか、という根本問題を未解決のまま残す。この意味で、量子状態の解釈の問題は、測定の解釈の問題に置き換えられただけである。また、量子状態を測定された観測可能量の固有状態へと還元する波動関数の収縮を仮定しても、物理的性質の「現実化」の問題が解決されるわけではないことにも注意しよう。そのような収縮は、それ自体では状態の存在論的地位を変えない。したがって、収縮後の系の状態が、収縮前より「物理的に実在する」と見なすことはできない。

量子状態の「実在論的」解釈の対極をなすのは、先に述べた、量子状態の純粋に認識論的な見解である。この解釈では、量子状態は私たちの知識だけを記述し、系の客観的な物理状態を記述するものではないと見なす（その例は文献\cite{ref044}、およびBallentineの「集団解釈」\cite{ref045,ref046}を参照）。この見解を経験的に反証するのは明らかに難しいが、それは、私たちの観測、したがって「知識」の基礎にある物理的実在という根本問題を扱うことを頑なに拒んでいる。

本書では、量子状態ベクトル（たとえば標準的なDirac記法でケット$\lvert\psi\rangle$と表す）は、個々の系の物理状態の完全な記述を与えるという、広く受け入れられている考えを採用する。このような量子状態の「完全性」を反映して、これらは通常、純粋と呼ばれる。それに対し、いわゆる混合状態は、純粋状態の古典的な集団にすぎず、2.4.2節で論じる。私たちの目標は、まだ発見されていない「隠れた変数」をその場しのぎに導入したり、私たちの「知識」を量子形式が記述する根本的な対象と見なしたりする安易な道を採る代わりに、量子形式の内部から私たちの観測の性質を説明することである。
